% Options for packages loaded elsewhere
\PassOptionsToPackage{unicode}{hyperref}
\PassOptionsToPackage{hyphens}{url}
\documentclass[
]{article}
\usepackage{xcolor}
\usepackage[margin=1in]{geometry}
\usepackage{amsmath,amssymb}
\setcounter{secnumdepth}{-\maxdimen} % remove section numbering
\usepackage{iftex}
\ifPDFTeX
  \usepackage[T1]{fontenc}
  \usepackage[utf8]{inputenc}
  \usepackage{textcomp} % provide euro and other symbols
\else % if luatex or xetex
  \usepackage{unicode-math} % this also loads fontspec
  \defaultfontfeatures{Scale=MatchLowercase}
  \defaultfontfeatures[\rmfamily]{Ligatures=TeX,Scale=1}
\fi
\usepackage{lmodern}
\ifPDFTeX\else
  % xetex/luatex font selection
\fi
% Use upquote if available, for straight quotes in verbatim environments
\IfFileExists{upquote.sty}{\usepackage{upquote}}{}
\IfFileExists{microtype.sty}{% use microtype if available
  \usepackage[]{microtype}
  \UseMicrotypeSet[protrusion]{basicmath} % disable protrusion for tt fonts
}{}
\makeatletter
\@ifundefined{KOMAClassName}{% if non-KOMA class
  \IfFileExists{parskip.sty}{%
    \usepackage{parskip}
  }{% else
    \setlength{\parindent}{0pt}
    \setlength{\parskip}{6pt plus 2pt minus 1pt}}
}{% if KOMA class
  \KOMAoptions{parskip=half}}
\makeatother
\usepackage{graphicx}
\makeatletter
\newsavebox\pandoc@box
\newcommand*\pandocbounded[1]{% scales image to fit in text height/width
  \sbox\pandoc@box{#1}%
  \Gscale@div\@tempa{\textheight}{\dimexpr\ht\pandoc@box+\dp\pandoc@box\relax}%
  \Gscale@div\@tempb{\linewidth}{\wd\pandoc@box}%
  \ifdim\@tempb\p@<\@tempa\p@\let\@tempa\@tempb\fi% select the smaller of both
  \ifdim\@tempa\p@<\p@\scalebox{\@tempa}{\usebox\pandoc@box}%
  \else\usebox{\pandoc@box}%
  \fi%
}
% Set default figure placement to htbp
\def\fps@figure{htbp}
\makeatother
\setlength{\emergencystretch}{3em} % prevent overfull lines
\providecommand{\tightlist}{%
  \setlength{\itemsep}{0pt}\setlength{\parskip}{0pt}}
\usepackage{graphicx}
\usepackage{float}
\usepackage{booktabs}
\usepackage{array}
\usepackage{bookmark}
\IfFileExists{xurl.sty}{\usepackage{xurl}}{} % add URL line breaks if available
\urlstyle{same}
\hypersetup{
  hidelinks,
  pdfcreator={LaTeX via pandoc}}

\author{}
\date{\vspace{-2.5em}}

\begin{document}

\section{Question}\label{question}

La gráfica y la tabla muestran parte de la información que recibe la
familia López en su factura telefónica del mes de enero.

\includegraphics[width=0.65\linewidth,height=\textheight,keepaspectratio]{grafico_consumo.png}

\includegraphics[width=10cm,height=\textheight,keepaspectratio]{tabla_estado_cuenta.png}

El tiempo adicional consumido por la familia López en enero fue:

\subsection{Answerlist}\label{answerlist}

\begin{itemize}
\tightlist
\item
  2 horas y 29 minutos.
\item
  6 horas y 13 minutos.
\item
  3 horas y 2 minutos.
\item
  2 horas y 8 minutos.
\end{itemize}

\section{Solution}\label{solution}

Para resolver este problema, debemos:

\begin{enumerate}
\def\labelenumi{\arabic{enumi}.}
\item
  \textbf{Identificar el plan actual}: Según la gráfica, el plan local
  actual es de \textbf{200 minutos}.
\item
  \textbf{Identificar el consumo de enero}: En el gráfico de barras, el
  consumo de enero (la última barra) es de \textbf{349 minutos}.
\item
  \textbf{Calcular el tiempo adicional}:
  \[\text{Tiempo adicional} = \text{Consumo enero} - \text{Plan actual}\]
  \[\text{Tiempo adicional} = 349 - 200 = 149 \text{ minutos}\]
\item
  \textbf{Convertir a horas y minutos}:

  \begin{itemize}
  \tightlist
  \item
    Horas: \(\lfloor 149 \div 60 \rfloor = 2\) hora(s)
  \item
    Minutos restantes: \(149 \mod 60 = 29\) minutos
  \end{itemize}
\end{enumerate}

\textbf{Respuesta correcta}: 2 horas y 29 minutos.

\subsection{Answerlist}\label{answerlist-1}

\begin{itemize}
\tightlist
\item
  Verdadero
\item
  Falso
\item
  Falso
\item
  Falso
\end{itemize}

\section{Meta-information}\label{meta-information}

exname: consumo\_telefonico\_adicional extype: schoice exsolution: 1000
exshuffle: TRUE exsection: Estadística/Interpretación de Gráficos/Nivel
2 extags: ICFES, Interpretación, Numérico-Variacional, Estadística,
Aplicado exdificultad: 2 excompetencia: Interpretación y Representación
excomponente: Numérico-Variacional expensamiento: Pensamiento Numérico
excontenido: Estadística (Genérico) exeje: Aplicado/Contextualizado

\end{document}
